\section{Contraste de hipótesis}

\begin{ejemplo}
    Sea $X\rightsquigarrow \cc{N}(\mu,1)$. Formular el problema de contrastar $\mu=15$ frente a $\mu=18$ y resolverlo de forma intuitiva usando un test no aleatorizado basado en $\ol{X}$ con tamaño $\alpha$.\\

    El problema consiste en contrastar si la media poblacional es igual a 15 frente a que sea igual a 18:
    \begin{align*}
        \begin{cases}
        H_0: \mu = 15 \\
        H_1: \mu = 18
        \end{cases}
    \end{align*}
    
    El espacio paramétrico considerado es, por tanto, $\Theta = \{15, 18\}$, dividido en $\Theta_0 = \{15\}$ y $\Theta_1 = \{18\}$. Sabemos que la media muestral $\overline{X}$ es un estadístico suficiente para $\mu$ y que su distribución es:
    \[
    \overline{X} \rightsquigarrow \mathcal{N}\left(\mu, \frac{1}{n}\right)
    \]
    \begin{itemize}
        \item Si $H_0$ es cierta ($\mu = 15$), los valores de $\overline{X}$ tenderán a situarse cerca de 15.
        \item Si $H_1$ es cierta ($\mu = 18$), los valores de $\overline{X}$ tenderán a concentrarse en torno a 18 (un valor mayor).
    \end{itemize}

    Por tanto, intuitivamente, la región de rechazo debe ser de cola a la derecha:
    \begin{equation*}
        \cc{C} = \left\{(x_1, \dots, x_n) \in \mathbb{R}^n : \overline{X} > k\right\}
    \end{equation*}
    donde $k$ es una constante crítica que queda determinada por el tamaño $\alpha$. Por tanto, el test no aleatorizado se define como:
    \begin{equation*}
        \varphi(x_1, \dots, x_n) =
        \begin{cases}
        1, & \text{si } \overline{X} > k \\
        0, & \text{si } \overline{X} \leq k
        \end{cases}
    \end{equation*}
    
    
    Dado que el test deberá tener tamaño $\alpha$, tenemos que:
    \begin{equation*}
        \alpha = P_{\mu=15}\left[\overline{X} > k\right] = 1 - P_{\mu=15}\left[\overline{X} \leq k\right] = 1 - F_{\overline{X}}(k; \mu=15)
    \end{equation*}

    Como $\overline{X} \rightsquigarrow \mathcal{N}\left(\mu, \nicefrac{1}{n}\right)$, tenemos que:
    \begin{equation*}
        1-\alpha = F_{\overline{X}}(k; \mu=15) = \Phi\left((k - 15)\sqrt{n}\right)
    \end{equation*}
    donde $\Phi$ es la función de distribución de la normal estándar $\mathcal{N}(0,1)$. Por tanto, dado $z_\alpha$ el valor de la normal estándar que deja a su derecha una probabilidad $\alpha$, es decir, $P[Z > z_\alpha] = \alpha$, tenemos que:
    \begin{equation*}
        (k-15)\sqrt{n} = z_\alpha \implies k = 15 + \frac{z_\alpha}{\sqrt{n}}
    \end{equation*}
\end{ejemplo}

\begin{ejemplo}
    Sea $(X_1,\ldots,X_5)$ una m.a.s. de $X\rightsquigarrow \cc{B}(1,p)$, con $p\in[0,1]$. Construir un test no aleatorizado para el contraste $H_0:p=\nicefrac{1}{2}$ frente a $H_1:p\neq\nicefrac{1}{2}$ con nivel de significación $\alpha = 0.1$.\\

    El problema consiste en contrastar si la probabilidad de éxito es igual a $\nicefrac{1}{2}$ frente a la alternativa bilateral:
    \begin{align*}
        \begin{cases}
        H_0: p = \nicefrac{1}{2} \\
        H_1: p \neq \nicefrac{1}{2}
        \end{cases}
    \end{align*}

    El espacio paramétrico es $\Theta = [0,1]$, particionado en los conjuntos $\Theta_0 = \{\nicefrac{1}{2}\}$ (hipótesis simple) y $\Theta_1 = [0,1]\setminus\{\nicefrac{1}{2}\}$ (hipótesis compuesta). Consideramos como estadístico del test el número total de éxitos en la muestra, $Y = \sum\limits_{i=1}^{5} X_i$, cuya distribución es $Y \rightsquigarrow \cc{B}(5,p)$.
    \begin{itemize}
        \item Si $H_0$ es cierta ($p = \nicefrac{1}{2}$), la distribución de $Y$ es simétrica respecto a su esperanza $E_{H_0}[Y] = 5 \cdot \nicefrac{1}{2} = 2.5$, por lo que las realizaciones de $Y$ tenderán a concentrarse en los valores centrales $\{2, 3\}$.
        \item Si $H_1$ es cierta ($p \neq \nicefrac{1}{2}$), la distribución se desplaza, concentrándose los valores hacia el $0$ (si $p < \nicefrac{1}{2}$) o hacia el $5$ (si $p > \nicefrac{1}{2}$).
    \end{itemize}

    Por tanto, intuitivamente, la región crítica o de rechazo $\cc{C}$ debe ser bilateral (de dos colas), rechazando cuando se observen valores extremos:
    \begin{equation*}
        \cc{C} = \left\{(x_1, \dots, x_5) \in \{0,1\}^5 : Y \leq k_1 \text{ o } Y \geq k_2\right\}
    \end{equation*}
    donde los puntos de corte verifican $0 \leq k_1 < k_2 \leq 5$. Por la simetría de la distribución binomial bajo $H_0$ ($p = \nicefrac{1}{2}$), repartimos el error de tipo I a partes iguales entre ambas colas, exigiendo que los valores críticos equidisten del centro $2.5$, es decir, $k_2 = 5 - k_1$.

    La condición para que el test tenga nivel de significación $\alpha = 0.1$ impone que la probabilidad de error de tipo I no supere dicha cota:
    \begin{equation*}
        P_{p=\nicefrac{1}{2}}[Y \in \cc{C}] = P_{p=\nicefrac{1}{2}}[Y \leq k_1] + P_{p=\nicefrac{1}{2}}[Y \geq 5 - k_1] \leq 0.1
    \end{equation*}
    Dado que $P_{p=\nicefrac{1}{2}}[Y = y] = \binom{5}{y}\left(\nicefrac{1}{2}\right)^5 = \frac{1}{32}\binom{5}{y}$ es simétrica respecto a $2.5$, ambas colas tienen igual probabilidad:
    \begin{equation*}
        P_{p=\nicefrac{1}{2}}[Y \leq k_1] = P_{p=\nicefrac{1}{2}}[Y \geq 5 - k_1] = \frac{1}{32}\sum_{i=0}^{k_1}\binom{5}{i}
    \end{equation*}
    Por tanto, la restricción sobre el nivel de significación se simplifica a:
    \begin{equation*}
        2 \cdot P_{p=\nicefrac{1}{2}}[Y \leq k_1] = \frac{2}{32}\sum_{i=0}^{k_1}\binom{5}{i} \leq 0.1 \iff \sum_{i=0}^{k_1}\binom{5}{i} \leq \frac{3.2}{2} = 1.6
    \end{equation*}

    Evaluamos sucesivamente los posibles valores enteros de $k_1 \in \{0, 1, 2\}$:
    \begin{itemize}
        \item Para $k_1 = 0$ (lo que implica $k_2 = 5 - 0 = 5$):
        \begin{equation*}
            \sum_{i=0}^{0} \binom{5}{i} = \binom{5}{0} = 1 \leq 1.6
        \end{equation*}
        El tamaño del test asociado es:
        \begin{equation*}
            P_{p=\nicefrac{1}{2}}[Y \in \{0, 5\}] = \frac{1}{32} + \frac{1}{32} = \frac{2}{32} = 0.0625 \leq 0.1
        \end{equation*}
        \item Para $k_1 = 1$ (lo que implica $k_2 = 5 - 1 = 4$):
        \begin{equation*}
            \sum_{i=0}^{1} \binom{5}{i} = \binom{5}{0} + \binom{5}{1} = 1 + 5 = 6 > 1.6
        \end{equation*}
        El tamaño ascendería a:
        \begin{equation*}
            P_{p=\nicefrac{1}{2}}[Y \in \{0, 1, 4, 5\}] = \frac{1 + 5 + 5 + 1}{32} = \frac{12}{32} = 0.375 > 0.1
        \end{equation*}

        En efecto, el tamaño del test superaría el nivel prefijado, por lo que no es admisible.
        \item Para $k_1 = 2$ (lo que implica $k_2 = 5 - 2 = 3$):
        \begin{equation*}
            \sum_{i=0}^{2} \binom{5}{i} = \binom{5}{0} + \binom{5}{1} + \binom{5}{2} = 1 + 5 + 10 = 16 > 1.6
        \end{equation*}
        El tamaño del test sería:
        \begin{equation*}
            P_{p=\nicefrac{1}{2}}[Y \in \{0, 1, 2, 3, 4, 5\}] = 1 > 0.1
        \end{equation*}
        Nuevamente, el tamaño del test superaría el nivel prefijado, por lo que no es admisible.
    \end{itemize}

    Por consiguiente, el test no aleatorizado óptimo (el de mayor potencia admisible sin superar el nivel $\alpha = 0.1$) se obtiene fijando $k_1 = 0$ y $k_2 = 5$. La región de rechazo resultante es:
    \begin{equation*}
        \cc{C} = \left\{(x_1, \dots, x_5) \in \{0,1\}^5 : \sum_{i=1}^5 x_i \in \{0, 5\}\right\}
    \end{equation*}
    cuyo tamaño efectivo es $0.0625$.
\end{ejemplo}

\begin{ejemplo}
    Sea $(X_1, \dots, X_5)$ una m.a.s. de $X\rightsquigarrow \cc{B}(1,p)$. Encontrar el test más potente de tamaño $\alpha = 0.05$ para resolver el problema de contraste:
    \begin{align*}
        \begin{cases}
        H_0: p = p_0 = \nicefrac{1}{2} \\
        H_1: p = p_1 = \nicefrac{3}{4}
        \end{cases}
    \end{align*}

    Sea $\theta_0=\nicefrac{1}{2}$ y $\theta_1=\nicefrac{3}{4}$. Notamos por $f_0\equiv f_{\theta_0}$ y $f_1\equiv f_{\theta_1}$ a las funciones de densidad de probabilidad de la muestra bajo $H_0$ y $H_1$, respectivamente. Consideramos el Test de Neyman-Pearson:
    \begin{equation*}
        \varphi(X_1, \dots, X_5) =
        \begin{cases}
            1, & \text{si } f_1(X_1, \dots, X_5) > k f_0(X_1, \dots, X_5) \\
            \gamma, & \text{si } f_1(X_1, \dots, X_5) = k f_0(X_1, \dots, X_5) \\
            0, & \text{si } f_1(X_1, \dots, X_5) < k f_0(X_1, \dots, X_5)
        \end{cases}
    \end{equation*}

    Buscamos por tanto calcular la función de densidad conjunta de la muestra:
    \begin{align*}
        f_p^n (x_1, \dots, x_n) &= \prod_{i=1}^{n} f_p(x_i) = \prod_{i=1}^{n} p^{x_i}(1-p)^{1-x_i} = p^{\sum\limits_{i=1}^{n} x_i}(1-p)^{n-\sum\limits_{i=1}^{n} x_i} \\
        &= p^{\sum\limits_{i=1}^{n} x_i}(1-p)^{n-\sum\limits_{i=1}^{n} x_i}>0\qquad \forall (x_1, \dots, x_n) \in \{0,1\}^n
    \end{align*}

    Notemos que, si $(x_1, \dots, x_n) \notin \{0,1\}^n$, entonces $f_p^n (x_1, \dots, x_n) = 0$ para cualquier $p\in[0,1]$. En cuyo caso se tiene que $f_1^n (x_1, \dots, x_n) = f_0^n (x_1, \dots, x_n) = 0$. Para $(x_1, \dots, x_n) \in \{0,1\}^n$, se define la razón de verosimilitudes:
    \begin{equation*}
        \lambda(x_1, \dots, x_n) = \frac{f_1^n(x_1, \dots, x_n)}{f_0^n(x_1, \dots, x_n)}
    \end{equation*}

    En tal caso, el Test de Neyman-Pearson se reduce a:
    \begin{equation*}
        \varphi(x_1, \dots, x_n) =
        \begin{cases}
            1, & \text{si } \lambda(x_1, \dots, x_n) > k \\
            \gamma, & \text{si } \lambda(x_1, \dots, x_n) = k \\
            0, & \text{si } \lambda(x_1, \dots, x_n) < k
        \end{cases}
    \end{equation*}
    
    
    
    Para $(x_1, \dots, x_n) \in \{0,1\}^n$, la razón de verosimilitudes es:
    \begin{align*}
        \lambda(x_1, \dots, x_n) &= \frac{f_1^n(x_1, \dots, x_n)}{f_0^n(x_1, \dots, x_n)} = \frac{\left(p_1\right)^{\sum\limits_{i=1}^{n} x_i}\left(1-p_1\right)^{n-\sum\limits_{i=1}^{n} x_i}}{\left(p_0\right)^{\sum\limits_{i=1}^{n} x_i}\left(1-p_0\right)^{n-\sum\limits_{i=1}^{n} x_i}} = \left(\frac{p_1}{p_0}\right)^{\sum\limits_{i=1}^{n} x_i}\left(\frac{1-p_1}{1-p_0}\right)^{n-\sum\limits_{i=1}^{n} x_i}
    \end{align*}

    Vemos que $\lambda(x_1, \dots, x_n)$ depende de la muestra únicamente a través del estadístico $T(X_1, \dots, X_n) = \sum\limits_{i=1}^{n} X_i$. Por tanto, $\lambda(x_1, \dots, x_n) = g(T(x_1, \dots, x_n))$, donde:
    \begin{align*}
        g(t) &= \left(\frac{p_1}{p_0}\right)^{t}\left(\frac{1-p_1}{1-p_0}\right)^{n-t}
    \end{align*}

    Para traducir las desigualdades del Test de Neyman-Pearson en términos del estadístico $T$, consideramos la función $g(t)$ y su comportamiento. Como $g(t)>0$ y el logaritmo es una función estrictamente creciente, podemos analizar la función logarítmica:
    \begin{align*}
        \ln g(t) &= t \ln\left(\frac{p_1}{p_0}\right) + (n-t) \ln\left(\frac{1-p_1}{1-p_0}\right)\\
        \dfrac{d}{dt} \ln g(t) &= \ln\left(\frac{p_1}{p_0}\right) - \ln\left(\frac{1-p_1}{1-p_0}\right) = \ln\left(\frac{p_1(1-p_0)}{p_0(1-p_1)}\right)>0
    \end{align*}
    donde en la última desigualdad se ha usado que $p_1>p_0$. Por tanto, $\ln g(t)$ es estrictamente creciente en $t$, y por ende, $g(t)$ también lo es. Esto implica que, para cada $k>0$, existe un único valor crítico $k'=g^{-1}(k)$. Para traducir las desigualdades del Test de Neyman-Pearson en términos del estadístico $T$, vemos gráficamente la Figura~\ref{fig:razon_verosimilitudes}.
    \begin{figure}[ht]
        \centering
        \begin{tikzpicture}[>=Stealth, scale=1.15]

            % --- Definición de estilos ---
            \tikzset{
                eje/.style={->, thick, black},
                guia/.style={dashed, thick, black!80},
                reg_sup/.style={ultra thick, orange!90!black},
                reg_inf/.style={ultra thick, red!75!black},
                punto_crit/.style={fill=cyan!80!blue, draw=black, thick}
            }

            % --- 1. Ejes de coordenadas ---
            \draw[eje] (-0.4, 0) -- (5.5, 0) node[right] {$T = \sum\limits_{i=1}^n X_i$};
            \draw[eje] (0, -0.4) -- (0, 4.3) node[above] {$\lambda = g(T)$};

            % --- 2. Curva creciente g(y) ---
            % Usamos una función exponencial suave g(x) = 0.45 * exp(0.42 * x)
            \draw[very thick, red!75!black, domain=0.2:3.2, samples=100] 
                plot (\x, {0.45 * exp(0.42 * \x)});
            \draw[very thick, orange!90!black, domain=3.2:4.8, samples=100] 
                plot (\x, {0.45 * exp(0.42 * \x)}) 
                node[above right, black] {$g(t)$ creciente};

            % --- 3. Coordenadas del punto de corte (k', k) ---
            % Para x = 3.2, y = 0.45 * exp(0.42 * 3.2) ≈ 1.725
            \def\kprima{3.2}
            \def\kval{1.725}

            % Líneas discontinuas de proyección
            \draw[guia] (0, \kval) -- (\kprima, \kval) -- (\kprima, 0);

            % Punto sobre la curva
            \filldraw[punto_crit] (\kprima, \kval) circle (2.2pt);

            % --- 4. Regiones coloreadas sobre los ejes ---
            % Eje vertical (lambda): lambda > k (arriba) y lambda < k (abajo)
            \draw[reg_sup] (0, \kval) -- (0, 3.9);
            \draw[reg_inf] (0, 0) -- (0, \kval);

            % Eje horizontal (T): T > k' (derecha) y T < k' (izquierda)
            \draw[reg_sup] (\kprima, 0) -- (5.1, 0);
            \draw[reg_inf] (0, 0) -- (\kprima, 0);

            % --- 5. Etiquetas de los puntos críticos ---
            \filldraw[punto_crit] (0, \kval) circle (2pt) node[left=4pt, black] {$k$};
            \filldraw[punto_crit] (\kprima, 0) circle (2pt) node[below=4pt, black] {$k'$};

            % --- 6. Etiquetas de las regiones (como en los apuntes) ---
            % Etiquetas en el eje vertical
            \node[left=4pt, orange!90!black, font=\small] at (0, 2.85) {$\lambda > k$};
            \node[left=4pt, red!75!black, font=\small] at (0, 0.85) {$\lambda < k$};

            % Etiquetas en el eje horizontal
            \node[below=4pt, orange!90!black, font=\small] at (4.2, 0) {$T > k'$};
            \node[below=4pt, red!75!black, font=\small] at (1.6, 0) {$T < k'$};

        \end{tikzpicture}
        \caption{Correspondencia entre la razón de verosimilitudes $\lambda$ y el estadístico $T$.}
        \label{fig:razon_verosimilitudes}
    \end{figure}
    
    Por tanto, el Test de Neyman-Pearson se puede reescribir en términos del estadístico $T$ como:
    \begin{equation*}
        \varphi(x_1, \dots, x_n) =
        \begin{cases}
            1, & \text{si } T(x_1, \dots, x_n) > k' \\
            \gamma, & \text{si } T(x_1, \dots, x_n) = k' \\
            0, & \text{si } T(x_1, \dots, x_n) < k'
        \end{cases}
    \end{equation*}

    Como se pide un test de tamaño $\alpha = 0.05$, debemos elegir $k'$ y $\gamma$ de manera que:
    \begin{align*}
        \alpha = 0.05 &= \sup_{p\in \Theta_0} E_{p}[\varphi(X_1, \dots, X_n)] = E_{p_0}[\varphi(X_1, \dots, X_n)] = \\&= P_{p_0}[T > k'] + \gamma P_{p_0}[T = k'] + \gamma \cancel{P_{p_0}[(X_1, \dots, X_n) \notin \{0,1\}^n]}
    \end{align*}

    Por tanto, despejando $\gamma$ obtenemos:
    \begin{align*}
        \gamma &= \frac{\alpha - P_{p_0}[T > k']}{P_{p_0}[T = k']}
    \end{align*}


    Sabiendo que $T \rightsquigarrow \cc{B}(n, p)$, buscamos los valores de $k'$ que nos permiten obtener $\gamma\in [0,1]$. Para $n=5$ y $p_0 = \nicefrac{1}{2}$, tenemos:
    \begin{align*}
        P_{p_0}[T > 5] &= 0 \\
        P_{p_0}[T > 4] &= P_{p_0}[T = 5] = \binom{5}{5}\left(\frac{1}{2}\right)^5 = \frac{1}{32} = 0.03125 \\
        P_{p_0}[T > 3] &= P_{p_0}[T = 4] + P_{p_0}[T = 5] = \binom{5}{4}\left(\frac{1}{2}\right)^5 + \binom{5}{5}\left(\frac{1}{2}\right)^5 = \frac{5 + 1}{32} = \frac{6}{32} = 0.1875
    \end{align*}

    Por tanto:
    \begin{itemize}
        \item Si $k' = 5$, entonces:
        \begin{equation*}
            \gamma = \frac{0.05 - 0}{P_{p_0}[T = 5]} = \frac{0.05}{0.03125} = 1.6 > 1
        \end{equation*}
        \item Si $k' = 4$, entonces:
        \begin{equation*}
            \gamma = \frac{0.05 - 0.03125}{P_{p_0}[T = 4]} = \frac{0.01875}{0.15625} = 0.12 \in [0,1]
        \end{equation*}
        \item Si $k' = 3$, entonces:
        \begin{equation*}
            \gamma = \frac{0.05 - 0.1875}{P_{p_0}[T = 3]} < 0
        \end{equation*}
    \end{itemize}

    Por tanto, el test más potente de tamaño $\alpha = 0.05$ se obtiene fijando $k' = 4$ y $\gamma = 0.12$. El test resultante es:
    \begin{equation*}
        \varphi(x_1, \dots, x_n) =
        \begin{cases}
            1, & \text{si } T(x_1, \dots, x_n) > 4 \\
            0.12, & \text{si } T(x_1, \dots, x_n) = 4 \\
            0, & \text{si } T(x_1, \dots, x_n) < 4
        \end{cases}
    \end{equation*}

    Por el Lema de Neyman-Pearson, este test es el más potente de tamaño $\alpha = 0.05$ para contrastar $H_0: p = \nicefrac{1}{2}$ frente a $H_1: p = \nicefrac{3}{4}$.
\end{ejemplo}


\begin{ejercicio} % // TODO: Hacer
    Sea $(X_1, \dots, X_n)$ una m.a.s. de $X$ cuya función de densidad viene dada por:
    \begin{equation*}
        f_\theta(x) = e^{\theta - x} \quad x \geq \theta
    \end{equation*}

    Encontrar el test más potente de tamaño $\alpha$ para resolver el problema de contraste:
    \begin{align*}
        \begin{cases}
        H_0: \theta = \theta_0 \\
        H_1: \theta = \theta_1
        \end{cases}
    \end{align*}
\end{ejercicio}